\section{Arm definitions}
\label{app:arms}

This appendix writes out each executed arm's exact update equation alongside the
five-slot taxonomic signature $\opcfg{R_{512}}{R_{2048}}{L_{512}}{L_{2048}}{\tau}$
of Section~\ref{sec:taxonomy} and Table~\ref{tab:arms}. Throughout, $W\in\mathbb
R^{m\times n}$ is a block matrix, $G_t$ its globally averaged accumulated
gradient, and
\[
M_t=\mu M_{t-1}+G_t,\qquad U_t=G_t+\mu M_t,\qquad \mu=0.95,\quad M_{-1}=0,
\]
the momentum and Nesterov source, with no $(1-\mu)$ factor and no bias
correction. Given a gradient-driven direction $D_t$ and a decay direction
$Q_t$,
\begin{equation}
\label{eq:update}
W_{t+1}=W_t-\eta_t\left(s_{m,n}D_t+\lambda Q_t\right),
\qquad s_{m,n}=\sqrt{\max(1,m/n)},\quad \lambda=0.1 .
\end{equation}
Both $D_t$ and $Q_t$ use pre-update weights and statistics. Factors are
$R=\In_n(C_t)$, $\hat R=\Out_n(C_t)$ from the input moment and
$L=\In_m(B_t)$, $\hat L=\Out_m(B_t)$ from the output moment, with $\In$ and
$\Out$ as in \eqref{eq:factors}; a superscript ``diag'' denotes the same
exponent applied elementwise to the moment's diagonal, normalized as
\[
a=\frac{d\,f}{\mathbf 1^{\Tr}f}+\epsilon\mathbf 1,
\qquad f=\diag(F),
\]
\[
\In^{\text{diag}}_d=\Diag\!\left(a^{-1/2}\right),\qquad
\Out^{\text{diag}}_d=\frac{\sqrt d\,\Diag\!\left(a^{-1/8}\right)}{\|a^{-1/8}\|_2}.
\]
Zero-trace statistics return identity factors. Powers in $\In$ and $\Out$ are
spectral; powers in the diagonal versions are elementwise.

\paragraph{Input-side MuonT (\muont{}), the baseline.}
For all seven projections,
\[
D_t=\Nten\!\left(U_tR\right)\hat R,\qquad Q_t=W_t\hat R^{\,2}.
\]
Both factors act on the right; decay is shaped by the squared outer input
factor.

\paragraph{Narrow arms, MLP down projections only.}
The other six projections take the baseline's update unchanged.
\begin{align*}
\texttt{outer\_left\_down}:\quad & D_t=\hat L\,\Nten\!\left(U_tR\right)\hat R, & Q_t&=W_t\hat R^{\,2},\\
\texttt{down\_hybrid}:\quad & D_t=\hat L\,\Nten\!\left(L\,U_t\right), & Q_t&=\hat L^{\,2}W_t,\\
\texttt{inner\_outer\_down}\ (\text{unexecuted}):\quad & D_t=\hat L\,\Nten\!\left(L\,U_tR\right)\hat R, & Q_t&=W_t\hat R^{\,2}.
\end{align*}
\texttt{outer\_left\_down} \emph{adds} a left factor after Newton--Schulz and
keeps every input factor; \texttt{down\_hybrid} \emph{replaces} the down
projection's input statistics, capture and roots with output ones, and uses
left-shaped decay. The 32 down matrices hold 33{,}554{,}432 weights, 12.6\% of
trainable parameters.

\paragraph{Broad arms, all 224 matrices.}
Each projection $p$ has its own output moment $B_p$; input moments remain
shared as in Table~\ref{tab:projections}.
\begin{align*}
\dense{}:\quad & D_{p,t}=\hat L_p\Nten\!\left(L_pU_{p,t}R_p\right)\hat R_p,\\
\leftall{}:\quad & D_{p,t}=\hat L_p\Nten\!\left(L_pU_{p,t}\right),\\
\texttt{outer\_left\_all}:\quad & D_{p,t}=\hat L_p\Nten\!\left(U_{p,t}R_p\right)\hat R_p,\\
\mlpdiag{}:\quad & D_{p,t}=\hat L^{\star}_p\Nten\!\left(L^{\star}_pU_{p,t}R_p\right)\hat R_p,
\end{align*}
with input-shaped decay $Q_{p,t}=W_{p,t}\hat R_p^{\,2}$ in every case except
\leftall{}, whose decay is left-shaped, $Q_{p,t}=\hat L_p^{\,2}W_{p,t}$.
Here, in \mlpdiag{}, the starred factors are dense for
$p\in\{q,k,v,o,\text{down}\}$ and diagonal for $p\in\{\text{gate},\text{up}\}$.
\alldiag{} is the same with starred factors diagonal everywhere;
\texttt{inner\_outer\_hidden} is the same with starred factors dense on the
width-512 group and \emph{absent} on $\text{gate},\text{up}$, whose cotangents
are then not observed at all. \texttt{inner\_outer\_all\_diag} replaces
\emph{both} sides by diagonals:
\[
D_{p,t}=\Lambda^{\text{out}}_{\text{out},p}\,
\Nten\!\left(\Lambda^{\text{out}}_{\text{in},p}\,U_{p,t}\,
\Lambda^{\text{in}}_{\text{in},p}\right)\Lambda^{\text{in}}_{\text{out},p},
\qquad
Q_{p,t}=W_{p,t}\left(\Lambda^{\text{in}}_{\text{out},p}\right)^{2},
\]
with the superscript naming the statistic's side and the subscript the inner or
outer root. This arm does not contain the baseline as a special case, which is
why its large regression does not isolate the output-side diagonal; that
isolation is what \alldiag{} provides.

\paragraph{Matched nulls.}
\nullall{} runs \dense{}'s path with $L_p=\hat L_p=I$ forced for all $p$;
\nullmlp{} does the same for \mlpdiag{}, preserving which groups would have been
dense and which diagonal. Both retain cotangent capture, pooled moments,
communication, root computation, host refresh, scheduling and checkpoint state.
Note that forcing identity left factors in \leftall{} would give
$D=\Nten(U)$, $Q=W$ --- a plain Muon-form update at the baseline's learning rate
rather than a baseline null --- which is why \leftall{}'s matched control is a
different arm, and why we did not run it: \leftall{}'s regression is large
enough ($+0.0087$) that no plausible path effect explains it.

\paragraph{Lagged arms.}
\texttt{inner\_outer\_all\_lag} and \texttt{inner\_outer\_mlpdiag\_lag} compute
exactly their parent's factors and apply them one update late,
\[
D_t=\hat L^{(t-1)}\,\Nten\!\left(L^{(t-1)}U_tR^{(t)}\right)\hat R^{(t)} ,
\]
with the input side current. The first update has no earlier rows and applies
input factors only. In the mixed arm both the dense and the diagonal left
factors advance together, so the arm has one left-factor age rather than two.

\paragraph{Fixed-age and fused arms.}
\outerdiag{} applies only the outer diagonal output factor,
$D_{p,t}=\hat L^{\text{diag}}_p\Nten(U_{p,t}R_p)\hat R_p$, and \fused{} is the
same equation with the post-Newton--Schulz products fused into three right
multiplications instead of four:
$-(\eta\lambda W\hat R+\eta s\hat L X)\hat R$. That is the same real arithmetic
in a different floating-point order; our CPU diagnostics bound the discrepancy
at a relative Frobenius error of $1.4\times10^{-7}$, and the two arms' paired
training-loss difference is $+0.000033$ with a standard error of $0.000049$ ---
indistinguishable from zero, which is the expected result for a change of
kernels rather than equations and a useful check on it. The fixed-age series
additionally pins input host-basis installation to the next update boundary, so
every logged installation has age exactly one and no refresh is discarded.

\paragraph{Execution notes that the equations do not capture.}
The baseline calls the upstream batched optimizer; every output-side arm uses a
per-matrix path with direct FP32 root products and separately formed direction
and decay. From the second wave onward the per-matrix path is replaced by a
batched one, \texttt{block\_update\_stack}, which issues one update per shape
group per rank instead of one per matrix, cutting kernel launches on the update
path from roughly $10^4$ to a few hundred per step; momentum matches the
per-matrix path bitwise and the applied delta differs only by FP32
reassociation inside the batched products, bounded at a relative $10^{-4}$.
Step times are therefore comparable within a wave and not across the first two.
